\documentclass[11pt]{article}
\usepackage[a4paper,margin=21mm]{geometry}
\usepackage{fontspec}
\setmainfont{Latin Modern Roman}
\newfontfamily\CJKfont{Noto Serif CJK SC}
\XeTeXlinebreaklocale "zh"
\XeTeXlinebreakskip=0pt plus 1pt
\usepackage{amsmath,amssymb,amsthm,amscd,graphicx,color}
\usepackage{xurl}
\usepackage[unicode,hidelinks]{hyperref}
\allowdisplaybreaks
\setlength\emergencystretch{3em}
\numberwithin{equation}{section}

\numberwithin{equation}{section}

\newtheorem{Theorem}{Theorem}[section]
\newtheorem{Corollary}[Theorem]{Corollary}
\newtheorem{Lemma}[Theorem]{Lemma}
\newtheorem{Proposition}[Theorem]{Proposition}
 { \theoremstyle{definition}
\newtheorem{Definition}[Theorem]{Definition}
\newtheorem{Example}[Theorem]{Example}
\newtheorem{Remark}[Theorem]{Remark} }

\usepackage{subcaption}
\usepackage{bm}
\usepackage{mathtools}

\newcommand{\Gabor}{{\mathcal G}}
\newcommand{\Fock}{\calF_0}
\newcommand{\wlogg}{\textit{w.l.o.g.}}
\newcommand{\nbh}{\Omega}
\newcommand{\ti}{\textit{t.i.}}
\newcommand{\ri}{\textit{r.i.}}
\newcommand{\lti}{\textit{u.t.i.}}
\newcommand{\bfR}{\mathbf{R}}
\newcommand{\bfconfig}{\mathbf{\Theta}}
\newcommand{\meso}{r}
\newcommand{\calA}{{\mathcal A}}
\newcommand{\calB}{{\mathcal B}}
\newcommand{\calC}{{\mathcal C}}
\newcommand{\calH}{{\mathcal H}}
\newcommand{\calI}{{\mathcal I}}
\newcommand{\calJ}{{\mathcal J}}
\newcommand{\calL}{{\mathcal L}}
\newcommand{\calK}{{\mathcal K}}
\newcommand{\calS}{{\mathcal S}}
\newcommand{\compset}{V}
\newcommand{\calP}{{\mathcal P}}
\newcommand{\calM}{{\mathcal M}}
\newcommand{\calV}{{\mathcal V}}
\newcommand{\calF}{{\mathcal F}}
\newcommand{\calR}{{\mathcal R}}
\newcommand{\calW}{{\mathcal W}}
\newcommand{\bfB}{{\mathbf B}}
\newcommand{\bfH}{{\mathbf H}}
\newcommand{\bfK}{{\mathbf K}}
\newcommand{\bfL}{{\mathbf L}}
\newcommand{\bfk}{{\mathbf k}}
\newcommand{\bfkappa}{{\bm \varkappa}}
\newcommand\modul{\tau_0}
\newcommand{\bulk}{\operatorname{bulk}}
\newcommand{\Int}{\operatorname{Int}}
\newcommand{\La}{L}
\newcommand{\Bz}{\hat{\bfB}}
\newcommand{\eqpot}{\check{Q}}
\newcommand{\config}{\Theta}
\newcommand{\trace}{\operatorname{trace}}
\newcommand{\erfc}{\operatorname{erfc}}
\newcommand{\erf}{\operatorname{erf}}
\newcommand{\Bulk}{\operatorname{bulk}}
\newcommand{\Boundary}{\operatorname{Boundary}}
\newcommand{\I}{{\mathrm{I}}}
\newcommand{\II}{{\mathrm{II}}}
\newcommand{\III}{{\mathrm{III}}}
\newcommand{\R}{{\mathbb R}}
\newcommand{\Z}{{\mathbb Z}}
\newcommand{\Q}{{\mathbb Q}}
\newcommand{\C}{{\mathbb C}}
\newcommand{\E}{{\mathbf E}}
\newcommand{\vro}{{\varrho}}
\newcommand{\fii}{{\varphi}}
\newcommand{\eps}{{\varepsilon}}
\newcommand{\diag}{\operatorname{diag}}
\newcommand{\re}{\operatorname{Re}}
\newcommand{\im}{\operatorname{Im}}
\newcommand{\ess}{\operatorname{ess}}
\newcommand{\Ham}{\mathbf{H}}
\newcommand{\calZ}{\mathcal{Z}}
\newcommand{\Pol}{\mathcal{W}}
\newcommand{\Poly}{\operatorname{Pol}}
\newcommand{\normalsp}{{\mathcal N}}
\newcommand{\normalpf}{{\mathcal Z}}
\newcommand{\normalvolfor}{{dM}}
\newcommand{\normalprob}{{\mathcal P}}
\newcommand{\Prob}{{\mathbf{P}}}
\newcommand{\volc}{{dV}}
\newcommand{\unitary}{{dU}}
\newcommand{\Tr}{\operatorname{Tr}}
\renewcommand{\d}{{\partial}}
\newcommand{\dbar}{\bar{\partial}}
\newcommand{\1}{\chi}
\newcommand{\Var}{\operatorname{Var}}
\newcommand{\Cov}{\operatorname{Cov}}
\newcommand{\lspan}{\operatorname{span}}
\newcommand{\dist}{\operatorname{dist}}
\newcommand{\supp}{\operatorname{supp}}
\newcommand{\diam}{\operatorname{diam}}
\newcommand{\Lap}{\Delta}
\newcommand{\lnorm}{\left\|}
\newcommand{\rnorm}{\right\|}
\newcommand{\drop}{S}
\newcommand{\Pf}{{\operatorname{Pf}}}
\newcommand{\sgn}{\operatorname{sgn}}

\def\bp{{\bar\partial}}
\def\wh{\widehat}
\def\norm#1{\lnorm {#1} \rnorm}
\def\lpar{\left (}
\def\rpar{\right )}
\def\labs{\left |}
\def\rabs{\right |}
\def\babs#1{\labs {#1} \rabs}
\def\alert#1{\textcolor{red}{#1}}
\newcommand{\RN}[1]{%
	\textup{\uppercase\expandafter{\romannumeral#1}}%
}
\newcommand{\J}{\mathbf{J}}
\newcommand{\D}{\mathbf{\mathcal{D}}}

\usepackage[skip=5pt,labelfont=bf,labelsep=period]{caption}

\makeatletter\newlabel{density}{{1.3}{}{Original source locator}{}{}}\makeatother
\makeatletter\newlabel{Fig. QGinibre}{{1}{}{Original source locator}{}{}}\makeatother
\makeatletter\newlabel{kappa bulk}{{2.6}{}{Original source locator}{}{}}\makeatother
\makeatletter\newlabel{Subsec_bdy Ginibre}{{3.2}{}{Original source locator}{}{}}\makeatother
\begin{document}
\begin{center}\Large Origin scaling kernels for symplectic Mittag-Leffler ensembles with reciprocal-integer radial exponent and point charge c>-1\end{center}
\noindent\textbf{Stable ID:} 1219\quad\textbf{Author:} Gernot Akemann; Sung-Soo Byun; Nam-Gyu Kang\par
\noindent\textbf{Work:} Scaling Limits of Planar Symplectic Ensembles\par
\noindent\textbf{Source:} \url{https://sigma-journal.com/2022/007/}\par
\noindent\textbf{Version:} Official SIGMA 18 (2022) article 007, DOI 10.3842/SIGMA.2022.007; journal PDF and native TeX acquired 2026-09-30, exact bytes pinned by SHA256\par
\noindent\textbf{Rights:} CC BY 4.0. Authors retain copyright. \url{https://creativecommons.org/licenses/by/4.0/}. Reuse and typesetting adaptation; original language and mathematical content preserved.\par
\noindent\textbf{Difficulty (prior selection preserved):} {\CJKfont H2: The proof derives a new Christoffel-Darboux differential system directly from the gamma-product double series of skew-orthogonal polynomials. It establishes the compact-uniform scaling kernel and solves the resulting integer-order inhomogeneous equation using a Mittag-Leffler fundamental system. An explicit Wronskian normalization and beta-integral coefficient identity remove every homogeneous contribution, yielding the full antisymmetric integral kernel.}\par
\medskip\noindent\textbf{Editorial boundary note (not author text):} Complete original Theorem2.4 at the singular origin for exactly reciprocal-integer lambda=1/m, every positive integer m and c>-1, including its m=1 and m>1 formulas and full definitions of F/g/W/G. The finite Gibbs/Pfaffian/radial skew-polynomial representation, exact micro-scale, branch/cocycle conventions, full current gamma-product limit and differential system, Wronskian/variation-of-parameters normalization and complete beta-integral coefficient proof are retained. General fractional Proposition4.1 is supporting core, not separately counted. The standard radial skew-polynomial identity, Mittag-Leffler homogeneous solutions and Abel/variation-of-parameters facts remain exact author references. The elementary first-order integrating-factor case is explicitly noted by the author rather than separately expanded. The last proof display retains its literal tilde-W after F extraction, alongside the earlier full integral and explicit tilde-W=(sign)F W identity; no symbol repair, editor proof or single-valued pre-kernel claim is added. Unselected Gaussian/Ward material and illustrative plots are excluded. Unused bbm import omitted after zero retained-call inventory.\par
\noindent\textbf{External original reference density:} Original macroscopic density convention, source215–221, standard cited equilibrium background; the selected singular micro-scale is explicit in2.13.\par
\noindent\textbf{External original reference Fig. QGinibre:} Original Figure1, an unselected sampling illustration mentioned only in rescaling commentary; no proof label or construction is located in that panel.\par
\noindent\textbf{External original reference kappa bulk:} Original Gaussian bulk reference kernel, source314–318, mentioned only as contrast in the selected Mittag-Leffler framing.\par
\noindent\textbf{External original reference Subsec\_bdy Ginibre:} Original Section3.2; its full cocycle definition is retained at native1089–1096 without importing the unselected Gaussian edge proof.\par
\medskip\hrule\medskip\noindent\textbf{Author text begins below.}\par
\setcounter{section}{1}\setcounter{equation}{0}
% Original source lines 164--172
Below we will consider more general potentials $Q$, where the joint probability distribution $\Prob_N$ of the point process $\boldsymbol{\zeta}=(\zeta_1,\dots,\zeta_N) \in \C^N$ is given by
\begin{gather}\label{Gibbs}
	{\rm d} \Prob _N(\boldsymbol{\zeta}):= \frac{1}{Z_N} {\rm e}^{-\bfH_N( \boldsymbol{\zeta} )} \prod_{j=1}^N {\rm d}A(\zeta_j).
\end{gather}
Here the Hamiltonian $\bfH_N$ is given by
\begin{gather} \label{Ham}
	\bfH_N( \boldsymbol{\zeta} ):=\sum_{j\neq k} \log \frac{1}{|\zeta_j-\zeta_k| \big|\zeta_j-\overline{\zeta}_k\big| }+\sum_{j=1}^{N} \left( \log \frac{1}{\big| \zeta_j- \overline{\zeta}_j \big|^2}+ NQ(\zeta_j) \right),
\end{gather}
the normalisation constant $Z_N$ is called the partition function, which turns $\Prob _N$ into a probability measure, and ${\rm d}A(\zeta):=\tfrac{1}{\pi}{\rm d}^2\zeta$ is the two-dimensional Lebesgue measure divided by~$\pi$.

\setcounter{section}{1}
% Original source lines 236--298
\section{Main results}\label{sec:main}
Let us now come to the main objects in this work. We denote by $\mathbb{D}(\eta,r)$ the disc with centre $\eta \in \C$ and radius $r$. For a given sequence of points $p_N$, the positive number $r_N=r_N(p_N)$ is called the \textit{micro-scale} if it satisfies
\begin{gather} \label{micro-scale}
	\int_{ \mathbb{D}(p_N,r_N) } \frac{ \Delta Q(\zeta) }{2} \,{\rm d}A(\zeta)= \frac{1}{N}.
\end{gather}
We drop the subscript and write $p \equiv p_N$ if the sequence does not depend on $N$.
By \eqref{density}, the micro-scale $r_N$ corresponds to the mean eigenvalue spacing in radial distance of the ensemble~\eqref{Gibbs} at the point $p\in S$
(cf.~\cite{BE} for a situation where $p$ is outside the droplet).
We define the rescaled process $\boldsymbol{z}=\{ z_j \}_{ j=1 }^N$ as follows: for all $j$,
\begin{gather} \label{rescaling}
	z_j:=
	\begin{cases}
		r_N^{-1} \cdot (\zeta_j-p) &\text{if } p \in \operatorname{int}(S) ,
		\\
		{\rm e}^{-{\rm i}\theta} r_N^{-1} \cdot (\zeta_j-p) &\text{if } p \in \partial S ,
	\end{cases}
\end{gather}
distinguishing the interior and the boundary of the droplet $S$.
Here the angle $\theta\equiv \theta_N \in \R$ is chosen so that ${\rm e}^{{\rm i}\theta}$ is outer normal to $\partial S$ at $p$, see Figure~\ref{Fig. QGinibre}(b)
inset
for an illustration of the rescaled process.
The $k$-point correlation function of
the rescaled process $\boldsymbol{z}$ is defined by
\begin{gather*}
R_{N,k}(z_1,\dots,z_k)
:= \lim_{\eps \downarrow 0} \frac{\mathbb{P}(\exists \text{ at least one particle in } \mathbb{D}(z_j,\eps),\, j=1,\dots,k ) }{\eps^{2k}}.
\end{gather*}
In \eqref{bfRNk def} a more standard definition of the $k$-point correlation function before rescaling denoted by $\bfR_{N,k}(\zeta_1,\dots,\zeta_k)$ is given.
Throughout the paper we distinguish these and all other objects (kernels, Hamiltonian, joint distribution) before rescaling by bold symbols.
For the precise relation on the level of the kernels see~\eqref{kappa_N}. It results into the following relation between the $k$-point correlation functions
\begin{gather*}
R_{N,k}(z_1,\dots, z_k)=r_N^{2k} \bfR_{N,k}(\zeta_1,\dots,\zeta_k).
\end{gather*}
It is well known that before and after rescaling the
set $\boldsymbol{z}$ forms a Pfaffian point process (see \cite{MR1928853,Mehta}), i.e., $R_{N,k}$ is expressed in terms of a certain $2\times 2$ matrix-valued kernel $K_N$ as
\begin{gather} \label{RNk Pf}
	R_{N,k}(z_1,\dots, z_k)=\prod_{j=1}^{k}
(\bar{z}_j-z_j) \Pf \big[ K_{N}(z_j,z_l) \big]_{ j,l=1 }^k ,
\end{gather}
where the kernel $K_N$ is of the form
\begin{gather*}
	K_N(z,w):={\rm e}^{ -\frac{N}{2} (Q( \zeta )+Q( \eta )) } \begin{pmatrix}
		\kappa_N(z,w) & \kappa_N(z,\bar{w}) \\
		\kappa_N(\bar{z},w) & \kappa_N(\bar{z},\bar{w})
	\end{pmatrix}.
\end{gather*}
Here $\Pf$ denotes a Pfaffian of the $2k \times 2k$ skew-symmetric matrix and
\begin{gather} \label{zeta eta}
\zeta=\begin{cases}
	p+r_N z &\text{if }p \in \textup{int}(S),
	\\
	p+{\rm e}^{{\rm i}\theta} r_N z &\text{if }p \in \partial S,
\end{cases} \qquad
 \eta=\begin{cases}
	 p+r_N w &\text{if }p \in \textup{int}(S),
	\\
	 p+{\rm e}^{{\rm i}\theta} r_N w &\text{if }p \in \partial S,
\end{cases}
\end{gather}
where $\theta$ is given as in \eqref{rescaling}.
The arguments \eqref{zeta eta} are given according to the rescaling \eqref{rescaling}.
For the spectral density at $k=1$, let us denote $R_N \equiv R_{N,1}$.


\setcounter{section}{2}\setcounter{Theorem}{3}\setcounter{equation}{11}
% Original source lines 460--523
Next, we investigate ensembles containing certain types of singularities at the origin. More precisely, we consider the \textit{Mittag-Leffler ensemble} whose associated potential $Q$ is of the form
\begin{gather} \label{Q ML}
	Q(\zeta)=|\zeta|^{2\lambda}-\frac{2c}{N} \log|\zeta|, \qquad \lambda>0, \quad c>-1.
\end{gather}
Here the condition $c>-1$ is required to guarantee $Z_N < \infty$, where $Z_N$ is the partition function in \eqref{Gibbs}. Note that when $c \not=0$, a ``conical singularity'' at the origin arises from an insertion of a point charge $c$.
On the other hand, since the limiting global density of the ensemble is $\Delta \tfrac12 Q(\zeta)=\tfrac{\lambda^2}2 |\zeta|^{2\lambda-2}$, when $\lambda>1$ (resp., $\lambda<1$) it vanishes (resp., diverges) at the origin.

We aim to discover the local statistics of the symplectic Mittag-Leffler ensemble, which provides ``non-standard'' or multi-critical universality classes \eqref{K ML} beyond \eqref{kappa bulk}.
Away from the origin we expect to be back in the Gaussian universality class of the Ginibre ensemble.
Note that the micro-scale at the origin is given here by
\begin{gather} \label{micro scale ML}
	r_N=\left(\frac{\lambda}{2}N\right)^{-\frac{1}{2\lambda}}.
\end{gather}
For each $\lambda>0$ and $c>-1$, the rescaled limiting local kernel $K_{\lambda,c}$ at $p=0$ is of the form
\begin{gather} \label{K ML}
	K_{\lambda,c}(z,w)
	={\rm e}^{ -\frac{ |z|^{2\lambda}+|w|^{2\lambda} }{\lambda} } \begin{pmatrix}
		\kappa_{\lambda,c}(z,w) & \kappa_{\lambda,c}(z,\bar{w}) \\
		\kappa_{\lambda,c}(\bar{z},w) & \kappa_{\lambda,c}(\bar{z},\bar{w})
	\end{pmatrix} ,
\end{gather}
where $\kappa_{\lambda,c}$ is a holomorphic function in both variables (with branch cuts).
Since the correlation functions do not depend on the specific choice of these branches, we will ignore this issue in the sequel, such a monodromy of the pre-kernel can also be interpreted as a cocycle. (Cf.\ Section~\ref{Subsec_bdy Ginibre} for the definition of a cocycle.)
In general, we derive a certain $(1/\lambda)$-order fractional differential equation for $\kappa_{\lambda,c}$ (Proposition~\ref{Thm_bulkS}), which can be recognised as a version of the Christoffel--Darboux formula for the kernel~\eqref{K ML}.
Moreover, for $\lambda=1/m$ $(m\in \mathbb{N})$, we obtain an explicit formula for $\kappa_{\lambda,c}$ by solving the associated differential equation of order~$m$.

To describe the local kernel of the Mittag-Leffer ensemble, let us first recall the definitions of the Mittag-Leffler functions. By definition, the \emph{two-parametric Mittag-Leffler function} $E_{a,b}$ is given by
\begin{gather} \label{ML function}
	E_{a,b}(z):=\sum_{k=0}^{\infty} \frac{z^k}{\Gamma(ak+b)},
\end{gather}
and the \emph{three-parametric $($Kilbas--Saigo$)$ Mittag-Leffler function} $E_{\alpha,m,l}$ is given by
\begin{gather} \label{ML3}
	E_{\alpha,m,l}(z):=1+\sum_{k=1}^\infty z^k \prod_{j=0}^{k-1} \frac{ \Gamma(\alpha(jm+l)+1) }{ \Gamma(\alpha(jm+l+1)+1) },
\end{gather}
see \cite[Chapters~4 and~5]{gorenflo2014mittag}.
For a positive integer $m$ we write
\begin{gather}\label{F integer}
	F_{m,c}(z):=z^{m(1+c)-1} E_{2m,m(1+c)}\big(z^{2m}\big).
\end{gather}
For each $j =1,\dots,m$, we write
\begin{gather} \label{m indep sol}
g_{j,m}(z):=z^{j-1} E_{m,2,1+\frac{j-1}{m}}\big(z^{2m}\big),
\end{gather}
and
\begin{gather} \label{W_j(z)}
W_{j,m}(z):=W( g_{1,m},\dots, g_{j-1,m}, g_{j+1,m}, \dots, g_{m,m} ),
\end{gather}
where $W$ is the Wronskian.

\begin{Theorem} \label{Thm_ML integer}
	For $Q(\zeta)=|\zeta|^{2/m}-\frac{2c}{N} \log|\zeta|$ with $\lambda=1/m$ $(m\in \mathbb{N})$, $c>-1$, we have
	\begin{gather*}
		\kappa_{\lambda,c}(z,w)=\left( \frac{2}{\lambda} \right)^{ \frac{1}{2\lambda} } (zw)^{\lambda-1} \widetilde{\kappa}_{\lambda,c} \left( \sqrt{ \frac{2}{\lambda} } z^\lambda, \sqrt{ \frac{2}{\lambda} } w^\lambda \right),
	\end{gather*}
	where for $m=1$,
	\begin{gather*}
	\widetilde{\kappa}_{1,c}(z,w)= \int_0^1 \big(z {\rm e}^{ \frac12 (1-s^2)z^2 } -w {\rm e}^{ \frac12 (1-s^2)w^2 }\big)F_{1,c}(szw)\,{\rm d}s,
	\end{gather*}
	and for $m>1$,
	\begin{gather*}
	\widetilde{\kappa}_{\lambda,c}(z,w)= \sum_{j=1}^{m} \frac{(-1)^{m-j}}{G(m+1)} \int_0^1 \big( z g_{j,m}(z)W_{j,m}(sz)-w g_{j,m}(w)W_{j,m}(sw)\big)F_{m,c}(szw)\,{\rm d}s.
	\end{gather*}
	Here $G(m+1):=\prod_{j=1}^{m-1} j!$ is the Barnes $G$-function.
\end{Theorem}

\setcounter{section}{3}\setcounter{subsection}{0}\setcounter{equation}{0}
% Original source lines 1004--1072
\subsection{Skew-orthogonal polynomials}
For a given potential $Q$, the \textit{anti-symmetric scalar product} $\langle \cdot | \cdot \rangle$ is given by
\begin{gather*}
	\langle h|g \rangle:= \int_\C \big( \bar{\zeta}-\zeta \big) {\rm e}^{-N Q(\zeta) } \big( h(\zeta) \overline{g(\zeta)}-\overline{h(\zeta)}g(\zeta) \big) \, {\rm d}A(\zeta).
\end{gather*}
By definition, a family $\{ q_k \}_{k=0}^\infty$ of polynomials is called
\textit{skew-orthogonal polynomials} associated with $Q$ if it satisfies
\begin{gather*}
	\langle q_{2k+1} | q_{2l} \rangle=-\langle q_{2l} | q_{2k+1} \rangle =s_k \, \delta_{kl}, \qquad \langle q_{2k+1} | q_{2l+1} \rangle=\langle q_{2l} | q_{2k} \rangle=0,
\end{gather*}
where $\delta_{kl}$ is the Kronecker delta and $s_k$ is a positive number that depends only on $k$.

For a radially symmetric potential $Q$, let
\begin{gather*}
	h_j:=\int_\C |\zeta|^{2j} {\rm e}^{-N Q(\zeta)} \, {\rm d}A(\zeta).
\end{gather*}
Then by \cite[Theorem 3.1]{akemann2021skew} (see also \cite[p.~7]{MR3066113})
\begin{gather} \label{skew op_rad}
	q_{2k+1}(\zeta)=\zeta^{2k+1}, \qquad q_{2k}(\zeta)=\zeta^{2k}+\sum_{l=0}^{k-1} \zeta^{2l} \prod_{j=0}^{k-l-1} \frac{h_{2l+2j+2} }{ h_{2l+2j+1} }
\end{gather}
form a family of skew-orthogonal polynomials associated with the potential $Q$. Here we have $s_k=2h_{2k+1}.$

Recall that the particle system \eqref{Gibbs} forms a Pfaffian point process, namely, the $k$-point correlation function
\begin{gather} \label{bfRNk def}
	\bfR_{N,k}(\zeta_1,\dots,\zeta_k):=\frac{1}{Z_N} \frac{N!}{(N-k)!} \int_{ \C^{N-k} } {\rm e}^{- \bfH_N( \boldsymbol{\zeta} ) } \prod_{j=k+1}^N {\rm d}A(\zeta_j)
\end{gather}
is expressed as
\begin{gather} \label{bfRNk Pf}
	\bfR_{N,k}(\zeta_1,\dots, \zeta_k)= \prod_{j=1}^{k} \big(\bar{\zeta}_j-\zeta_j\big) \Pf \big[ \bfK_{N}(\zeta_j,\zeta_l) \big]_{ j,l=1 }^k,
\end{gather}
where the $2 \times 2$ matrix-valued kernel $\bfK_N$ is of the form
\begin{gather} \label{bfK Pf}
	\bfK_N(\zeta,\eta)={\rm e}^{ -NQ(\zeta)/2-NQ(\eta)/2 }
	\begin{pmatrix}
		\bfkappa_N(\zeta,\eta) & \bfkappa_N(\zeta,\bar{\eta}) \\
		\bfkappa_N(\bar{\zeta},\eta) & \bfkappa_N(\bar{\zeta},\bar{\eta})
	\end{pmatrix}.
\end{gather}
In particular, we write $\bfR_{N} \equiv \bfR_{N,1}$ for the $1$-point function.
It is well known that the pre-kernel~$\bfkappa_N$ takes the form
\begin{gather} \label{bfkappaN}
	\bfkappa_N(\zeta,\eta)=\sum_{k=0}^{N-1} \frac{q_{2k+1}(\zeta) q_{2k}(\eta) -q_{2k}(\zeta) q_{2k+1}(\eta)}{s_k},
\end{gather}
see \cite{MR1928853}, where the quaternionic determinant $Q\det$ is used instead of the Pfaffian $\Pf$. We also refer to \cite{Mehta} for the relation between $Q\det$ and $\Pf$.

The relation between the pre-kernels $\bfkappa_N$ and $\kappa_N$
for the change of variables \eqref{zeta eta}
at point $p\in S$ is given as
\begin{gather} \label{kappa_N}
	\kappa_N(z,w):=r_N^3
	\begin{cases}
		\bfkappa_N( p+r_Nz, p+r_Nw) &\text{if } p \in \textup{int}(S),
		\\
	\bfkappa_N\big( p+r_Nz {\rm e}^{{\rm i}\theta}, p+r_Nw {\rm e}^{{\rm i}\theta} \big)&\text{if } p \in \partial S.
	\end{cases} 	
\end{gather}
Thus the rescaled process $\boldsymbol{z}$ retains its Pfaffian structure in terms of the rescaled matrix-valued kernel
\begin{gather} \label{K_N}
	K_N(z,w):={\rm e}^{ -\frac{N}{2} (Q( \zeta )+Q( \eta )) } \begin{pmatrix}
		\kappa_N(z,w) & \kappa_N(z,\bar{w}) \\
		\kappa_N(\bar{z},w) & \kappa_N(\bar{z},\bar{w})
	\end{pmatrix},
\end{gather}
cf.~\eqref{zeta eta} for the unscaled variables $\zeta$ and $\eta$ inside the arguments of the potential~$Q$. We remark that in the factor $r_N^3$ in~\eqref{kappa_N}, one factor $r_N$ comes from the term $\prod_{j=1}^k \big(\bar{\zeta}_j-\zeta_j\big)$ in~\eqref{bfRNk Pf} (see~\eqref{RNk Pf}), whereas the other two factors originate from the change of variables for the pre-kernel. This leads precisely to~\eqref{RNk Pf} as follows
	\begin{gather*}
		 R_{N,k}(z_1,\dots, z_k)
		 = 	r_N^{2k} \prod_{j=1}^{k} \big(\bar{\zeta}_j-\zeta_j\big) \Pf \big[ \bfK_{N}(\zeta_j,\zeta_l) \big]_{ j,l=1 }^k
		= \prod_{j=1}^{k} \frac{\bar{\zeta}_j-\zeta_j}{r_N} \Pf \big[ K_{N}(z_j,z_l) \big]_{ j,l=1 }^k.
	\end{gather*}


% Original source lines 1089--1096
By definition, a function $c(z,w)$ is called a \textit{cocycle} if there exists a (continuous) unimodular function $h$ (i.e., $h(z)=h(1/\bar{z})$) such that $c(z,w) = h(z)h(w).$ Since cocycles cancel out when multiplying the pre-kernel and taking the Pfaffian, if $(c_N)_{N=1}^\infty$ is a sequence of cocycles, then~$c_N\kappa_N$ is an equivalent realisation of the pre-kernel~$\kappa_N$
for the same point process $\boldsymbol{z}$. Let us denote by~$c_N \cdot K_N$ the correlation kernel associated with $c_N \kappa_N$, i.e.,
\begin{gather*}
c_N \cdot K_N(z,w):= {\rm e}^{ -\frac{N}{2} (Q( \zeta )+Q( \eta )) } \begin{pmatrix}
	 c_N(z,w)\, \kappa_N(z,w) & c_N(z,\bar{w})\, \kappa_N(z,\bar{w}) \\
	c_N(\bar{z},w)\, \kappa_N(\bar{z},w) & c_N(\bar{z},\bar{w}) \, \kappa_N(\bar{z},\bar{w})
	\end{pmatrix}.
\end{gather*}

\setcounter{section}{3}
% Original source lines 1291--1545
\section{Scaling limit of the Mittag-Leffler ensemble}

In this section, we construct a fractional differential equation satisfied by the local kernel of the Mittag-Leffler ensemble with potentials $Q(\zeta)=|\zeta|^{2\lambda}-(2c/N)\log|\zeta|$, $\lambda>0$ and $c>-1$ (Proposition~\ref{Thm_bulkS}). As a consequence, we prove Theorem~\ref{Thm_ML integer}.

\subsection{Christoffel--Darboux type formula}
For the Mittag-Leffler potentials \eqref{Q ML}, the orthogonal norm $h_k$ is given by
\begin{gather*}
h_k=\int_\C |\zeta|^{2k} {\rm e}^{-N Q(\zeta)} \, {\rm d}A(\zeta)=\frac{\Gamma\big( \frac{k+c+1}{\lambda} \big)}{\lambda N^{\frac{k+c+1}{\lambda}} } .
\end{gather*}
Then by \eqref{skew op_rad}, we have
\begin{gather} \label{SOP ML}
	q_{2k+1}(\zeta)=\zeta^{2k+1}, \qquad q_{2k}(\zeta)=\zeta^{2k}+ \sum_{l=0}^{k-1} \frac{\zeta^{2l}}{N^{\frac{k-l}{\lambda}}} \prod_{j=0}^{k-l-1} \frac{ \Gamma\big( \frac{2l+2j+3+c}{\lambda} \big) }{ \Gamma\big( \frac{2l+2j+2+c}{\lambda} \big) } ,
\end{gather}
and
\begin{gather} \label{sk ML}
	s_k=2h_{2k+1}=\frac{2}{\lambda } \frac{\Gamma( \frac{2k+c+2}{\lambda} )}{N^{\frac{2k+c+2}{\lambda}} } .
\end{gather}
To describe the Christoffel--Darboux type formula, let us recall that the Caputo fractional derivative $D^\nu_z$ is given by
\begin{gather*}
	D^\nu_z f(z):=\frac{1}{\Gamma(n-\nu)} \int_{0}^{z} (z-u)^{n-\nu-1} f^{(n)}(u) \, {\rm d}u, \qquad n-1 < \nu <n, \quad n \in \mathbb{N},
\end{gather*}
see \cite[Section 2.4]{MR2218073}. Here $f^{(n)}$ is the usual $n$:th derivative.

Recall that $p=0$
%%
in the rescaling \eqref{zeta eta} here,
and $K_{\lambda,c}$ is given by \eqref{K ML}.

\begin{Proposition} \label{Thm_bulkS}
	For each $\lambda>0$ and $c>-1$, there exists a sequence of cocycles $(c_N)_{N=1}^\infty$ such that
	\begin{gather*}
			\lim_{N \to \infty} c_N\cdot K_N(z,w)=K_{\lambda,c}(z,w)
			={\rm e}^{ -\frac{ |z|^{2\lambda}+|w|^{2\lambda} }{\lambda} } \begin{pmatrix}
				\kappa_{\lambda,c}(z,w) & \kappa_{\lambda,c}(z,\bar{w}) \\
				\kappa_{\lambda,c}(\bar{z},w) & \kappa_{\lambda,c}(\bar{z},\bar{w})
			\end{pmatrix}
		\end{gather*}
	uniformly for $z$, $w$ in compact subsets of~$\C$, where
	\begin{gather*}
		\kappa_{\lambda,c}(z,w)=\left( \frac{2}{\lambda} \right)^{ \frac{1}{2\lambda} } (zw)^{\lambda-1} \widetilde{\kappa}_{\lambda,c} \left( \sqrt{ \frac{2}{\lambda} } z^\lambda, \sqrt{ \frac{2}{\lambda} } w^\lambda \right).
	\end{gather*}
	Here $\widetilde{\kappa}_{\lambda,c}$ is given in terms of a function $\widetilde{G}$ as
	\begin{gather*}
	\widetilde{\kappa}_{\lambda,c}(z,w)=\widetilde{G}(z,w)-\widetilde{G}(w,z)
	\end{gather*}
 and the following fractional differential equations hold:
	\begin{gather} \label{G zw FDE}
		D_z^{ \frac{1}{\lambda} } \widetilde{G}(z,w)= z^{ \frac{1}{\lambda} } \widetilde{G}(z,w)+(zw)^{\frac{1+c}{\lambda}-1} E_{\frac{2}{\lambda} \frac{1+c}{\lambda} } \big( (zw)^{\frac{2}{\lambda}} \big),
\\
			D_z^{ \frac{1}{\lambda} } \widetilde{\kappa}_{\lambda,c}(z,w)
			= z^{ \frac{1}{\lambda} } \widetilde{\kappa}_{\lambda,c}(z,w)+ (zw)^{ \frac{1+c}{\lambda}-1 } E_{ \frac{1}{\lambda},\frac{1+c}{\lambda} } \big( (zw)^{\frac{1}{\lambda}} \big)
			\nonumber\\
\hphantom{D_z^{ \frac{1}{\lambda} } \widetilde{\kappa}_{\lambda,c}(z,w)=}{}
 -\frac{ \Gamma\big(\frac{1+c}{\lambda}\big) }{\Gamma\big( \frac{c}{\lambda}\big) \Gamma\big( \frac{2+c}{\lambda}\big) } (zw)^{ \frac{1+c}{\lambda}-1 }(w/z)^{\frac{1}{\lambda}} E_{\frac{1}{\lambda},2,3+c-\lambda}\big(w^{\frac{2}{\lambda}}\big) .\label{FDE_bulkS}
		\end{gather}
\end{Proposition}


We call the equation \eqref{FDE_bulkS} (generalised) Christoffel--Darboux formula.
Such a differential equation satisfied by the correlation kernel, that makes the asymptotic analysis possible, is broadly called Christoffel--Darboux type formula, see, e.g., \cite{MR1917675,byun2021lemniscate,MR3450566}.

We also remark that the inhomogeneous term $(zw)^{ \frac{1+c}{\lambda}-1 } E_{ \frac{1}{\lambda},\frac{1+c}{\lambda} } \big( (zw)^{\frac{1}{\lambda}} \big)$ corresponds to the kernel of the complex Mittag-Leffler ensemble, see~\cite{ameur2018random}.
Such a relation has been observed for other models as well, which include the elliptic Ginibre ensemble~\cite{BE} and its chiral counter part~\cite{MR2180006}.
See also~\cite{MR1762659,MR1675356} for a similar relation for Hermitian ensembles, which gives an expression of the kernel of the symplectic ensemble in terms of a small number of orthogonal polynomials.

\begin{proof}[Proof of Proposition~\ref{Thm_bulkS}]
	Let us define
	\begin{gather*} %\label{kappa tilde0 bS}
		\widetilde{G}_{N}(z,w):=\sum_{k=0}^{N-1}z^{ \frac{2k+2+c}{\lambda}-1 } \prod_{j=0}^k \frac{ \Gamma\big( \frac{2j+1+c}{\lambda} \big) }{\Gamma\big( \frac{2j+2+c}{\lambda} \big)}
		\sum_{l=0}^{k} \frac{ w^{ \frac{2l+1+c}{\lambda}-1 } }{\Gamma\big( \frac{2l+1+c}{\lambda}\big)} \prod_{j=1}^l \frac{\Gamma\big( \frac{2j+c}{\lambda} \big)}{ \Gamma\big(\frac{2j-1+c}{\lambda}\big) } .
	\end{gather*}
	Here and henceforth we use the convention that if $l=0$, then $\prod_{j=1}^l \varphi=1$.
	By \eqref{SOP ML}, \eqref{sk ML} and~\eqref{micro scale ML}, the kernel $K_N$ is given by
	\begin{gather*}
			K_N(z,w)
			={\rm e}^{ -\frac{ |z|^{2\lambda}+|w|^{2\lambda} }{\lambda} } \begin{pmatrix}
				\kappa_N(z,w) & \kappa_N(z,\bar{w}) \\
				\kappa_N(\bar{z},w) & \kappa_N(\bar{z},\bar{w})
			\end{pmatrix},
		\end{gather*}
	where
	\begin{gather*}
		\kappa_N(z,w):=G_N(z,w)-G_N(w,z), \\ G_N(z,w):=\left( \frac{2}{\lambda} \right)^{\frac{1}{2\lambda}} (zw)^{\lambda-1} \widetilde{G}_N \left( \sqrt{ \frac{2}{\lambda} } z^\lambda, \sqrt{ \frac{2}{\lambda} } w^\lambda \right).
	\end{gather*}
	
	By definition of the Caputo derivative, we have that if $k>
	\lceil \nu \rceil$
	\begin{gather} \label{Caputo monomial}
		D_z^\nu z^k =\frac{\Gamma(k+1)}{\Gamma(k-\nu+1)}z^{k-\nu}
	\end{gather}
	and vanishes otherwise.
	Let us write $\widetilde{G}:=\lim\limits_{N \to \infty} \widetilde{G}_N,$ where the convergence is uniform on compact subsets of $\C$. The existence of the limit follows from the Weierstrass $M$-test.
	By direct computations using \eqref{Caputo monomial}, we have
	\begin{gather}
\widetilde{G}(z,w)-\frac{1}{\Gamma\big(\frac{2+c}{\lambda}\big)} z^{ \frac{2+c}{\lambda}-1 } w^{ \frac{1+c}{\lambda}-1 }
			\nonumber\\
\qquad{} =\sum_{k=1}^{\infty} z^{ \frac{2k+2+c}{\lambda}-1 } \prod_{j=0}^k \frac{ \Gamma\big( \frac{2j+1+c}{\lambda} \big) }{\Gamma\big( \frac{2j+2+c}{\lambda} \big)}
			\sum_{l=0}^{k} \frac{ w^{ \frac{2l+1+c}{\lambda}-1 } }{\Gamma\big( \frac{2l+1+c}{\lambda}\big)} \prod_{j=1}^l \frac{\Gamma\big( \frac{2j+c}{\lambda} \big)}{ \Gamma\big(\frac{2j-1+c}{\lambda}\big) } .\label{G zw til lam c 0}
		\end{gather}
	Note in particular that as $z\to 0,$
	\begin{gather} \label{G zw til asym}
		\widetilde{G}(z,w) \sim \frac{w^{ \frac{1+c}{\lambda}-1 } }{\Gamma\big(\frac{2+c}{\lambda}\big)} z^{ \frac{2+c}{\lambda}-1 }.
	\end{gather}
	Applying the operator $D_z^{ \frac{1}{\lambda} }$ to the identity \eqref{G zw til lam c 0}, we obtain
	\begin{gather*}
 D_z^{ \frac{1}{\lambda} } \widetilde{G}(z,w)-\frac{ 1}{\Gamma\big( \frac{1+c}{\lambda} \big)}(zw)^{ \frac{1+c}{\lambda}-1 }
		\\
\qquad{} = \sum_{k=1}^{\infty} z^{ \frac{2k+1+c}{\lambda}-1 } \prod_{j=1}^{k}\frac{ \Gamma \big( \frac{2j-1+c}{\lambda} \big) }{\Gamma\big( \frac{2j+c}{\lambda} \big)} \sum_{l=0}^{k} \frac{ w^{ \frac{2l+1+c}{\lambda}-1 } }{\Gamma\big( \frac{2l+1+c}{\lambda} \big)} \prod_{j=1}^l \frac{\Gamma\big(\frac{2j+c}{\lambda}\big) }{ \Gamma\big( \frac{2j-1+c}{\lambda} \big) }
		\\
\qquad{} = z^{ \frac{1}{\lambda} } \sum_{k=0}^{\infty} z^{ \frac{2k+2+c}{\lambda}-1 } \prod_{j=0}^k \frac{ \Gamma \big( \frac{2j+1+c}{\lambda} \big) }{\Gamma\big( \frac{2j+2+c}{\lambda} \big)} \sum_{l=0}^{k+1} \frac{ w^{ \frac{2l+1+c}{\lambda}-1 } }{\Gamma\big( \frac{2l+1+c}{\lambda} \big)} \prod_{j=1}^l \frac{\Gamma\big(\frac{2j+c}{\lambda}\big) }{ \Gamma\big( \frac{2j-1+c}{\lambda} \big) } ,
	\end{gather*}
	which leads to
	\begin{gather*}
		D_z^{ \frac{1}{\lambda} } \widetilde{G}(z,w)= z^{ \frac{1}{\lambda} } \widetilde{G}(z,w)
		+\sum_{k=0}^{\infty} \frac{ (zw)^{ \frac{2k+1+c}{\lambda}-1 } }{\Gamma\big( \frac{2k+1+c}{\lambda} \big)}.
	\end{gather*}
	Now \eqref{G zw FDE} follows from \eqref{ML function}.
	
	Similarly, by rearranging the terms, we have
	\begin{gather}
 \widetilde{G}(w,z)-\frac{ 1 }{\Gamma\big( \frac{1+c}{\lambda}\big)} z^{ \frac{1+c}{\lambda}-1 } \sum_{k=0}^{\infty} w^{ \frac{2k+2+c}{\lambda}-1 } \prod_{j=0}^k \frac{ \Gamma\big( \frac{2j+1+c}{\lambda} \big) }{\Gamma\big( \frac{2j+2+c}{\lambda} \big)}
			\nonumber\\
\qquad{} =\sum_{k=1}^{\infty} w^{ \frac{2k+2+c}{\lambda}-1 } \prod_{j=0}^k \frac{ \Gamma\big( \frac{2j+1+c}{\lambda} \big) }{\Gamma\big( \frac{2j+2+c}{\lambda} \big)} \sum_{l=1}^{k} \frac{ z^{ \frac{2l+1+c}{\lambda}-1 } }{\Gamma\big( \frac{2l+1+c}{\lambda} \big)} \prod_{j=1}^l \frac{\Gamma\big(\frac{2j+c}{\lambda}\big) }{ \Gamma\big( \frac{2j-1+c}{\lambda} \big) } .\label{G wz til lam c 0}
		\end{gather}
	Then by \eqref{ML3}, we have that as $z\to 0$,
	\begin{gather*} %\label{G wz til asym}
		\widetilde{G}(w,z)\sim \frac{1 }{ \Gamma\big( \frac{2+c}{\lambda}\big) } w^{ \frac{2+c}{\lambda}-1 } E_{\frac{1}{\lambda},2,3+c-\lambda}\big(w^{\frac{2}{\lambda}}\big) z^{ \frac{1+c}{\lambda}-1 }
	\end{gather*}
	Again, by applying the operator $D_z^{ \frac{1}{\lambda} }$ to \eqref{G wz til lam c 0}, we obtain
	\begin{gather*}
D_z^{ \frac{1}{\lambda} } \widetilde{G}(w,z)-\frac{ \Gamma\big(\frac{1+c}{\lambda}\big) }{\Gamma\big( \frac{c}{\lambda}\big) \Gamma\big( \frac{2+c}{\lambda}\big) } (zw)^{ \frac{1+c}{\lambda}-1 }(w/z)^{\frac{1}{\lambda}} E_{\frac{1}{\lambda},2,3+c-\lambda}\big(w^{\frac{2}{\lambda}}\big)\\
\qquad{} = z^{ \frac{1}{\lambda} } \sum_{k=1}^{\infty} w^{ \frac{2k+2+c}{\lambda}-1 } \prod_{j=0}^k \frac{ \Gamma\big( \frac{2j+1+c}{\lambda} \big) }{\Gamma\big( \frac{2j+2+c}{\lambda} \big)} \sum_{l=0}^{k-1} \frac{ z^{ \frac{2l+1+c}{\lambda}-1 } }{\Gamma\big( \frac{2l+1+c}{\lambda} \big)} \prod_{j=1}^l \frac{\Gamma\big(\frac{2j+c}{\lambda}\big) }{ \Gamma\big( \frac{2j-1+c}{\lambda} \big) }
		\\
\qquad{} = z^{ \frac{1}{\lambda} } \widetilde{G}(w,z)- \sum_{k=0}^{\infty} \frac{(zw)^{ \frac{2k+2+c}{\lambda}-1 }}{\Gamma\big( \frac{2k+2+c}{\lambda} \big)}.
	\end{gather*}
	Then by \eqref{ML function}, we have
	\begin{gather}
D_z^{ \frac{1}{\lambda} } \widetilde{G}(w,z)= z^{ \frac{1}{\lambda} } \widetilde{G}(w,z) -(zw)^{\frac{2+c}{\lambda}-1} E_{\frac{2}{\lambda},\frac{2+c}{\lambda}} \big((zw)^{\frac{2}{\lambda}}\big) \nonumber\\
\hphantom{D_z^{ \frac{1}{\lambda} } \widetilde{G}(w,z)=}{}
+\frac{ \Gamma(\frac{1+c}{\lambda}) }{\Gamma( \frac{c}{\lambda}) \Gamma( \frac{2+c}{\lambda}) } (zw)^{ \frac{1+c}{\lambda}-1 }(w/z)^{\frac{1}{\lambda}} E_{\frac{1}{\lambda},2,3+c-\lambda}\big(w^{\frac{2}{\lambda}}\big) . \label{G wz FDE}
		\end{gather}
	Now \eqref{FDE_bulkS} follows from \eqref{G zw FDE}, \eqref{G wz FDE} and \eqref{ML function}.
\end{proof}


\subsection[Bulk singularities of the order lambda=1/m with m in N]{Bulk singularities of the order $\boldsymbol{\lambda=\frac1m}$ with $\boldsymbol{m\in \mathbb{N}}$}

A point $p$ in which the eigenvalue density $\frac12\Delta Q(p)$ vanishes or diverges is called bulk singularity, see \cite{ameur2018random} and references therein. In this subsection, we shall consider the case $\lambda=\frac1m$, where $m$ is a positive integer and prove Theorem~\ref{Thm_ML integer} for general $c>-1$.

\begin{proof}[Proof of Theorem~\ref{Thm_ML integer}]
	By \eqref{G zw FDE} and \eqref{G zw til asym}, the function $\widetilde{G}$ is a unique solution to the $m$-th order linear (ordinary) differential equation
	\begin{gather} \label{G zw ODE}
		\partial_z ^m \widetilde{G}(z,w)=z^m \widetilde{G}(z,w)+F_{m,c}(zw),
	\end{gather}
	which satisfies the asymptotic behaviour
	\begin{gather} \label{G zw asym int}
		\widetilde{G}(z,w) \sim \frac{w^{m+mc-1}}{ \Gamma(2m+mc) } z^{2m+mc-1},\qquad z \to 0.
	\end{gather}
	Here the function $F_{m,c}$ is given by \eqref{F integer}. Note that the $m$ initial conditions are provided by differentiating \eqref{G zw asym int}.
	
We aim to solve the initial value problem \eqref{G zw ODE} and~\eqref{G zw asym int}.
	It is well known that the functions~$g_{j,m}$ given by \eqref{m indep sol} provide the solutions to the homogeneous equation of~\eqref{G zw ODE}, see, e.g., \cite[Theorem~5.18]{gorenflo2014mittag}.
	By~\eqref{ML3}, we have that
	\begin{gather} \label{gj asymp}
		g_{j,m}(z) \sim z^{j-1} ,\qquad z \to 0.
	\end{gather}
	Using this, it is easy to observe that
	\begin{gather*}
		W(g_{1,m},\dots,g_{m,m})(z)|_{z=0}=\prod_{j=0}^{m-1} j!=G(m+1),
	\end{gather*}
	where $G$ is the Barnes $G$-function \cite[Section~5.17]{olver2010nist}. Due to Abel's identity, for $m>1$, this leads to
	\begin{gather} \label{Wronskian g1m}
		W(g_{1,m},\dots,g_{m,m})(z) \equiv G(m+1).
	\end{gather}
	On the other hand for $m=1$, we have obviously $W(g_1)(z)=g_1(z).$
	
	For each $j \in \{1,\dots, m\}$, let us write
	\begin{gather*}
		\widetilde{W}_{j,m}(z,w):=\det
		\begin{bmatrix}
			g_{1,m}(z) & \dots & g_{j-1,m}(z) & 0 & g_{j+1,m}(z) & \dots & g_{m,m}(z)
			\\
			g_{1,m}'(z) & \dots & g_{j-1,m}'(z) & 0 & g_{j+1,m}'(z) & \dots & g_{m,m}'(z)
			\\
			\vdots & \vdots & \vdots & \vdots & \vdots & \vdots & \vdots
			\\
			g_{1,m}^{(m-1)}(z) & \dots & g_{j-1,m}^{(m-1)}(z) & F_{m,c}(zw) & g_{j+1,m}^{(m-1)}(z) & \dots & g_{m,m}^{(m-1)}(z)
		\end{bmatrix}.
	\end{gather*}
	Note here that by \eqref{W_j(z)}, we have
	\begin{gather} \label{Wj zw F Wj}
		\widetilde{W}_{j,m}(z,w)=(-1)^{m-j} F_{m,c}(zw) W_{j,m}(z).
	\end{gather}
	
	From now on, we shall only consider the case $m>1$. The other case $m=1$ follows along the same lines. In this case, one can also easily solve the associated linear non-homogenous ODE of degree $1$ employing an integrating factor.
	
	By virtue of the method of variation of parameters and \eqref{Wronskian g1m}, one can write
	\begin{align*}
			\widetilde{G}(z,w)&= \int_0^z \frac{1}{W(g_{1,m},\dots,g_{m,m})(s)}\sum_{j=1}^{m} 	\widetilde{W}_{j,m}(s,w) g_{j,m}(z)\,{\rm d}s+ \sum_{j=1}^{m} C_{j,m}(w) g_{j,m}(z)
			\\
			&=\frac{1}{G(m+1)}\int_0^z \sum_{j=1}^{m} 	\widetilde{W}_{j,m}(s,w) g_{j,m}(z)\,{\rm d}s+ \sum_{j=1}^{m} C_{j,m}(w) g_{j,m}(z)
	\end{align*}
	for some functions $C_{j,m}$. We shall show that in the above expression $\widetilde{G}$ satisfies the asymptotic behaviour~\eqref{G zw asym int} if and only if $C_{j,m} \equiv 0$ for every~$j$.
	
	By \eqref{gj asymp} and \eqref{Wronskian g1m}, direct computations of the determinant give that as $z \to 0$,
	\begin{gather*}
	W_{j,m}(z) \sim \frac{G(m+1)}{(m-1)!} \binom{m-1}{j-1} z^{m-j}.
	\end{gather*}
	On the other hand, by \eqref{ML function}, as $z \to 0$,
	\begin{gather*}
	F_{m,c}(zw) \sim \frac{w^{m+mc-1}}{\Gamma(m+mc)} z^{m+mc-1}.
	\end{gather*}
	Combining these two equations, we have that as $z\to 0$,
	\begin{gather*}
	\widetilde{W}_{j,m}(z,w) \sim \frac{w^{m+mc-1}}{\Gamma(m+mc)} (-1)^{m-j} \frac{G(m+1)}{(m-1)!} \binom{m-1}{j-1} z^{2m+mc-1-j}.
	\end{gather*}
	This leads to
	\begin{gather*}
	\frac{1}{G(m+1)}\int_0^z \widetilde{W}_{j,m}(s,w) g_{j,m}(z)\,{\rm d}s\\
\qquad{} \sim \frac{w^{m+mc-1}}{\Gamma(m+mc)} \frac{1}{(m-1)!} \binom{m-1}{j-1} \frac{(-1)^{m-j} }{2m+mc-j} z^{2m+mc-1}.
	\end{gather*}
	
	Next, we show the identity
	\begin{gather}\label{Comb identity}
		\frac{1}{(m-1)!} \sum_{j=1}^{m} \binom{m-1}{j-1} \frac{(-1)^{m-j} }{2m+mc-j} = \frac{\Gamma(m+mc)}{\Gamma(2m+mc)}.
	\end{gather}
	By the binomial theorem, we have
	\begin{gather*}
	\sum_{k=0}^{m-1} \binom{m-1}{k} x^{m+mc-1+k} = x^{m+mc-1}(1+x)^{m-1}.
	\end{gather*}
	Integrating this identity, we have
	\begin{gather*}
	\sum_{k=0}^{m-1} \binom{m-1}{k} \frac{x^{m+mc+k}}{m+mc+k} = \int_0^x t^{m+mc-1}(1+t)^{m-1}\,{\rm d}t.
	\end{gather*}
	Letting $x=-1$ and using the change of the variable $t=-s$, we have
	\begin{align*}
		\sum_{k=0}^{m-1} \binom{m-1}{k} \frac{(-1)^{k}}{m+mc+k} &= \int_0^1 s^{m+mc-1}(1-s)^{m-1}\,{\rm d}s
		\\
		&=B(m+mc,m)=(m-1)! \frac{\Gamma(m+mc)}{\Gamma(2m+mc)},
	\end{align*}
	where $B$ is the beta function. This gives~\eqref{Comb identity}.
	
	Using \eqref{Comb identity}, we have
	\begin{gather*}
	\frac{1}{G(m+1)}\int_0^z \sum_{j=1}^{m} 	\widetilde{W}_{j,m}(s,w) g_{j,m}(z)\,{\rm d}s \sim \frac{w^{m+mc-1}}{ \Gamma(2m+mc) } z^{2m+mc-1}.
	\end{gather*}
	Thus by \eqref{G zw asym int}, we obtain that $C_{j,m} \equiv 0$ for all $j \in \{1,\dots, m\}$.
	Therefore by~\eqref{Wj zw F Wj}, we conclude that
	\begin{align*}
		\widetilde{G}(z,w)&= \frac{1}{G(m+1)}\int_0^1 \sum_{j=1}^{m} 	\widetilde{W}_{j,m}(sz,w) z g_{j,m}(z)\,{\rm d}s
		\\
		&= \frac{1}{G(m+1)}\int_0^1 \sum_{j=1}^{m} (-1)^{m-j} 	\widetilde{W}_{j,m}(sz) z g_{j,m}(z)F_{m,c}(szw)\,{\rm d}s.
	\end{align*}
	This completes the proof.
\end{proof}
\begin{thebibliography}{99}
\bibitem[1]{MR1762659}
Adler M., Forrester P.J., Nagao T., van Moerbeke P., Classical skew orthogonal
 polynomials and random matrices, \href{https://doi.org/10.1023/A:1018644606835}{\textit{J.~Statist. Phys.}} \textbf{99}
 (2000), 141--170, \href{https://arxiv.org/abs/solv-int/9907001}{arXiv:solv-int/9907001}.

\bibitem[2]{MR2180006}
Akemann G., The complex {L}aguerre symplectic ensemble of non-{H}ermitian
 matrices, \href{https://doi.org/10.1016/j.nuclphysb.2005.09.039}{\textit{Nuclear Phys.~B}} \textbf{730} (2005), 253--299,
 \href{https://arxiv.org/abs/hep-th/0507156}{arXiv:hep-th/0507156}.

\bibitem[6]{akemann2021skew}
Akemann G., Ebke M., Parra I., Skew-orthogonal polynomials in the complex plane
 and their {B}ergman-like kernels, \href{https://doi.org/10.1007/s00220-021-04230-8}{\textit{Comm. Math. Phys.}} \textbf{389}
 (2022), 621--659, \href{https://arxiv.org/abs/2103.12114}{arXiv:2103.12114}.

\bibitem[15]{ameur2018random}
Ameur Y., Kang N.-G., Seo S.-M., The random normal matrix model: insertion of
 a~point charge, \href{https://doi.org/10.1007/s11118-021-09942-z}{\textit{Potential Anal.}}, {t}o appear, \href{https://arxiv.org/abs/1804.08587}{arXiv:1804.08587}.

\bibitem[19]{MR1917675}
Bertola M., Eynard B., Harnad J., Duality, biorthogonal polynomials and
 multi-matrix models, \href{https://doi.org/10.1007/s002200200663}{\textit{Comm. Math. Phys.}} \textbf{229} (2002), 73--120,
 \href{https://arxiv.org/abs/nlin.SI/0108049}{arXiv:nlin.SI/0108049}.

\bibitem[21]{BE}
Byun S.-S., Ebke M., Universal scaling limits of the symplectic elliptic
 {G}inibre ensemble, \href{https://arxiv.org/abs/2108.05541}{arXiv:2108.05541}.

\bibitem[23]{byun2021lemniscate}
Byun S.-S., Lee S.-Y., Yang M., Lemniscate ensembles with spectral singularity,
 \href{https://arxiv.org/abs/2107.07221}{arXiv:2107.07221}.

\bibitem[32]{gorenflo2014mittag}
Gorenflo R., Kilbas A.A., Mainardi F., Rogosin S.V., Mittag-{L}effler
 functions, related topics and applications, \textit{Springer Monographs in
 Mathematics}, \href{https://doi.org/10.1007/978-3-662-43930-2}{Springer}, Heidelberg, 2014.

\bibitem[34]{MR3066113}
Ipsen J.R., Products of independent quaternion {G}inibre matrices and their
 correlation functions, \href{https://doi.org/10.1088/1751-8113/46/26/265201}{\textit{J.~Phys.~A: Math. Theor.}} \textbf{46} (2013),
 265201, 16~pages, \href{https://arxiv.org/abs/1301.3343}{arXiv:1301.3343}.

\bibitem[35]{MR1928853}
Kanzieper E., Eigenvalue correlations in non-{H}ermitean symplectic random
 matrices, \href{https://doi.org/10.1088/0305-4470/35/31/308}{\textit{J.~Phys.~A: Math. Gen.}} \textbf{35} (2002), 6631--6644,
 \href{https://arxiv.org/abs/cond-mat/0109287}{arXiv:cond-mat/0109287}.

\bibitem[39]{MR2218073}
Kilbas A.A., Srivastava H.M., Trujillo J.J., Theory and applications of
 fractional differential equations, \textit{North-Holland Mathematics
 Studies}, Vol.~204, Elsevier Science~B.V., Amsterdam, 2006.

\bibitem[41]{MR3450566}
Lee S.-Y., Riser R., Fine asymptotic behavior for eigenvalues of random normal
 matrices: ellipse case, \href{https://doi.org/10.1063/1.4939973}{\textit{J.~Math. Phys.}} \textbf{57} (2016), 023302,
 29~pages, \href{https://arxiv.org/abs/1501.02781}{arXiv:1501.02781}.

\bibitem[44]{Mehta}
Mehta M.L., Random matrices, 2nd ed., Academic Press, Inc., Boston, MA, 1991.

\bibitem[45]{olver2010nist}
Olver F.W.J., Lozier D.W., Boisvert R.F., Clark C.W. (Editors), N{IST} handbook
 of mathematical functions, Cambridge University Press, Cambridge, 2010.

\bibitem[48]{MR1675356}
Widom H., On the relation between orthogonal, symplectic and unitary matrix
 ensembles, \href{https://doi.org/10.1023/A:1004516918143}{\textit{J.~Statist. Phys.}} \textbf{94} (1999), 347--363,
 \href{https://arxiv.org/abs/solv-int/9804005}{arXiv:solv-int/9804005}.

\end{thebibliography}
\end{document}
